\chapter{Schistosomiasis}
\index{Schistosomiasis}
\label{sec:Schistosomiasis}

\section{Overview}
Schistosomiasis (bilharziasis) is a trematode infection caused by blood flukes of the genus \textit{Schistosoma}: \textit{S. haematobium} (urinary), \textit{S. mansoni} (intestinal), \textit{S. japonicum}, \textit{S. intercalatum}, and \textit{S. mekongi} (intestinal), with 200 million infected and 800 million at risk, primarily in sub-Saharan Africa, with transmission through contact with fresh water containing infected snails.
\section{Causes}
This condition happens because of cercariae penetrating human skin, transforming into schistosomula, migrating to the portal venous system, maturing into adult worms, and producing eggs that become trapped in tissues, causing granulomatous inflammation, scarring, organomegaly, and obstruction.     
\section{Risk Factors}

\begin{itemize}[noitemsep]
  \item Genetic predisposition and family history
  \item Advancing age
  \item Lifestyle factors (diet, exercise, smoking, alcohol)
  \item Environmental exposures
  \item Underlying medical conditions
  \item Medication use
\end{itemize}
\section{Signs and Symptoms}

\subsection{Common symptoms}
\begin{itemize}[noitemsep]
  \item Pain or discomfort in the affected area
  \end{itemize}

\subsection{Emergency warning signs}
\begin{itemize}[noitemsep]
  \item Severe or worsening symptoms of schistosomiasis require immediate medical evaluation
\end{itemize}
\section{Progression and Stages}
The condition may progress through different stages of severity over time. Early stages often involve mild or subtle symptoms that may be overlooked. Moderate stages involve more noticeable symptoms affecting daily function. Advanced stages may involve significant impairment and complications.
\section{Complications}
You may experience acute schistosomiasis (Katayama fever): fever, cough, urticaria, muscle pain, and hepatosplenomegaly 2--6 weeks after exposure

\begin{itemize}[noitemsep]
  \item Chronic \textit{S. haematobium}: hematuria, dysuria, hydronephrosis, and bladder scarring, with squamous cell carcinoma as a long-term complication
  \item Chronic \textit{S. mansoni} and \textit{S. japonicum}: hepatosplenomegaly, periportal scarring (Symmers' pipe-stem scarring), portal high blood pressure, and esophageal varices.
  \item Disease progression leading to worsening symptoms
  \item Secondary infections or injuries
  \item Side effects of treatment
  \item Reduced quality of life
  \item Emotional and psychological impact
\end{itemize}
\section{Diagnosis}
Doctors diagnose this using microscopic identification of eggs in stool or urine. Comprehensive medical history and physical examination Laboratory tests to assess organ function and rule out other conditions Imaging studies to visualize affected structures Specialized testing depending on the affected system Consultation with relevant specialists Monitoring of disease progression over time

\begin{itemize}[noitemsep]
  \item Comprehensive medical history and physical examination
  \item Laboratory tests to assess organ function
  \item Imaging studies to visualize affected structures
  \item Specialized testing depending on the affected system
  \item Consultation with relevant specialists
\end{itemize}
\section{Prevention}
\begin{itemize}[noitemsep]
  \item Maintain a healthy lifestyle with balanced diet and exercise
  \item Avoid known risk factors
  \item Attend regular medical check-ups for early detection
  \item Follow recommended screening guidelines
\end{itemize}
\section{Treatment Options}
Treatment options include praziquantel. Medications to manage symptoms and address underlying mechanisms Lifestyle modifications and self-care measures Physical therapy and rehabilitation Surgical intervention when indicated Regular monitoring and follow-up care Patient education and support services

\begin{itemize}[noitemsep]
  \item Medications to manage symptoms and address underlying mechanisms
  \item Lifestyle modifications and self-care measures
  \item Physical therapy and rehabilitation
  \item Surgical intervention when indicated
  \item Regular monitoring and follow-up care
\end{itemize}
\section{Living with It}
\begin{itemize}[noitemsep]
  \item Follow your treatment plan as prescribed
  \item Maintain regular follow-up with your healthcare provider
  \item Make necessary lifestyle adjustments
  \item Seek support from family, friends, or support groups
  \item Stay informed about your condition
\end{itemize}
\section{Outlook}
The outlook is good with treatment. The outlook varies depending on the specific condition, severity at diagnosis, and response to treatment. Early intervention generally leads to better outcomes. Many conditions can be effectively managed with appropriate treatment. Regular follow-up care is important for monitoring and treatment adjustment.
